% ==============================================================================
% Phụ lục B: Thẩm định Toàn diện 20 Tiêu chí PROBAST+AI
% Tham chiếu: Collins et al. / Moons et al., The BMJ 2024 & Annals of Internal Medicine 2025
% Tác giả phụ trách: Thành viên 2 (Biomedical Data Engineer & Privacy) - Backlog 10 Nhiệm vụ 2.2
% ==============================================================================
\section{Phụ Lục: Báo Cáo Thẩm Định Chi Tiết 20 Câu Hỏi Tín Hiệu PROBAST+AI}
\label{sec:appendix_probast_ai}

Công cụ **PROBAST+AI** (\textit{Prediction model Risk Of Bias ASsessment Tool for Artificial Intelligence}) \cite{probastai2024bmj,moons2025probast} được thiết kế chuyên biệt để đánh giá nguy cơ sai lệch (\textit{Risk of Bias}) và mối lo ngại về tính áp dụng (\textit{Applicability Concerns}) của các mô hình học máy trong y tế. 

Bảng~\ref{tab:probast_20_questions} trình bày chi tiết kết quả tự kiểm toán độc lập trên toàn bộ 20 câu hỏi tín hiệu thuộc 4 miền phương pháp luận, tích hợp đầy đủ các kết quả phân tích thống kê chuyên sâu từ Giai đoạn 3 (Backlog 08 và Backlog 09).

\begin{table}[htbp]
\centering
\caption{Đánh giá chi tiết 20 câu hỏi tín hiệu theo chuẩn PROBAST+AI cho hệ thống M-CDSS.}
\label{tab:probast_20_questions}
\tiny
\begin{tabularx}{\textwidth}{@{}cp{4.5cm}cp{5.2cm}@{}}
\toprule
\textbf{STT} & \textbf{Câu hỏi tín hiệu (Signaling Question)} & \textbf{Đánh giá} & \textbf{Bằng chứng kỹ thuật \& Biện luận khoa học trong M-CDSS} \\
\midrule
\multicolumn{4}{@{}l}{\textbf{MIỀN 1: NGƯỜI THAM GIA (PARTICIPANTS) --- Nguy cơ: CAO (HIGH) | Tính áp dụng: CAO (HIGH CONCERN)}} \\
1.1 & Nguồn dữ liệu có phù hợp và đại diện cho bối cảnh lâm sàng dự kiến áp dụng không? & Chưa rõ & Dữ liệu bán tổng hợp từ Synthea EHR kết hợp MTSamples, chưa phản ánh đầy đủ mô hình bệnh tật và văn phong bệnh án thực tế tại Việt Nam. \\
1.2 & Tiêu chí thu nhận và loại trừ người tham gia có phù hợp không? & Có & Thu nhận toàn bộ ca bệnh người lớn có triệu chứng hô hấp/tuần hoàn; không loại trừ tùy tiện ca bệnh nặng. \\
1.3 & Nguy cơ sai lệch lựa chọn mẫu (Selection Bias) có được kiểm soát không? & Có & Cân bằng 4 nhóm bệnh mục tiêu ($25\%$ mỗi nhóm, 500 ca/nhóm) trên quy mô toàn đoàn hệ. \\
\midrule
\multicolumn{4}{@{}l}{\textbf{MIỀN 2: BIẾN DỰ BÁO (PREDICTORS) --- Nguy cơ: THẤP/CHƯA RÕ | Tính áp dụng: CHƯA RÕ (UNCLEAR)}} \\
2.1 & Các biến dự báo có được định nghĩa và thu thập tương đồng cho mọi người tham gia không? & Có & Hợp đồng dữ liệu cố định chuẩn hóa 13 biến số bảng và văn bản tự do theo cùng một quy cách. \\
2.2 & Quá trình đánh giá biến dự báo có được thực hiện độc lập với kết cục không? & Có & Biến số được ghi nhận lúc tiếp nhận khám, độc lập với nhãn chẩn đoán kết luận sau cùng. \\
2.3 & Biến dự báo có sẵn sàng tại thời điểm mô hình dự kiến được áp dụng không? & Có & Toàn bộ 13 biến số và triệu chứng cơ năng đều có sẵn ngay tại phòng khám ban đầu. \\
2.4 & Có loại bỏ hoàn toàn các biến số vi phạm chuỗi thời gian hoặc lộ nhãn trực tiếp không? & Có & Module \texttt{TemporalAligner} khóa thời gian $t < t_{\text{diag}}$; danh sách đen cấm xét nghiệm tiêu chuẩn vàng ($FEV_1/FVC$, Siêu âm $EF$, Cấy đờm vi sinh) và che từ khóa chẩn đoán \texttt{[REDACTED\_DIAGNOSIS]}. \\
\midrule
\multicolumn{4}{@{}l}{\textbf{MIỀN 3: KẾT CỤC CHẨN ĐOÁN (OUTCOME) --- Nguy cơ: CAO (HIGH) | Tính áp dụng: CAO (HIGH CONCERN)}} \\
3.1 & Kết cục chẩn đoán có được định nghĩa theo tiêu chuẩn quốc tế công nhận không? & Có & Đối chiếu chuẩn với GINA 2024, GOLD 2024, ESC 2023 và Bộ Y tế Việt Nam. \\
3.2 & Các biến dự báo có được LOẠI BỎ hoàn toàn khỏi định nghĩa kết cục không? & Có & Không sử dụng định nghĩa lặp vòng; khắc phục sai sót của các nghiên cứu tiền nhiệm bị JMIR cảnh báo. \\
3.3 & Kết cục có được xác định tương đồng cho mọi người tham gia không? & Có & Toàn bộ hồ sơ được gán nhãn theo cùng danh mục phân loại mã bệnh ICD-10. \\
3.4 & Người xác định kết cục có được làm mù đối với các biến dự báo tiềm năng không? & Có & Nhãn mã bệnh được gán độc lập trong quy trình hành chính trước khi nghiên cứu AI được tiến hành. \\
3.5 & Khoảng thời gian giữa biến dự báo và kết cục có hợp lý và nhất quán không? & Có & Biến dự báo thu thập lúc tiếp nhận ($t_0$), chẩn đoán xác lập trong cùng đợt khám ($t_{\text{diag}}$). \\
\midrule
\multicolumn{4}{@{}l}{\textbf{MIỀN 4: PHÂN TÍCH SỐ LIỆU (ANALYSIS) --- Nguy cơ: THẤP (LOW) | Tính áp dụng: THẤP (LOW CONCERN)}} \\
4.1 & Cỡ mẫu và số lượng ca biến cố (Events Per Variable - EPV) có đủ lớn không? & Có & $N = 2.000$ với 500 biến cố/nhóm, tỷ lệ $\text{EPV} \ge 38$, vượt xa ngưỡng khuyến nghị tối thiểu $\ge 10-20$. \\
4.2 & Các biến số liên tục có được xử lý thích hợp mà không bị phân nhóm tùy tiện không? & Có & Giữ nguyên giá trị số thực liên tục của sinh hiệu và xét nghiệm, chuẩn hóa qua Layer Normalization. \\
4.3 & Tất cả đối tượng tham gia ban đầu có được giữ lại trong phân tích không? & Có & Không xóa dòng chứa dữ liệu thiếu (\textit{No Complete Case drop}); 100\% đối tượng được giữ lại. \\
4.4 & Dữ liệu khuyết thiếu (Missing Data) có được xử lý minh bạch và đúng phương pháp không? & Có & Áp dụng thuật toán MICE chỉ fit trên Train và transform sang Val/Test, kèm 2 biến chỉ báo thiếu (\texttt{wbc\_missing}, \texttt{crp\_missing}). \\
4.5 & Có tránh việc chọn lọc biến dựa trên phân tích đơn biến (Univariable Screening) không? & Có & Chọn lọc đặc trưng dựa trên tri thức lâm sàng tiên nghiệm kết hợp Dropout và Weight Decay. \\
4.6 & Cấu trúc phân cụm bệnh nhân có được kiểm soát khi chia tập dữ liệu không? & Có & Stratified Group Splitting theo \texttt{patient\_id} và \texttt{source\_note\_group}; kiểm toán $0\%$ rò rỉ giữa Train, Val và Test. \\
4.7 & Các chỉ số đo lường hiệu năng có phản ánh đầy đủ năng lực phân biệt và độ tin cậy không? & Có & Báo cáo đa chiều: AUROC, AUPRC, Macro-F1, ECE, Brier Score kèm khoảng tin cậy 95\% Bootstrap (2.000 lần lặp) và hiệu chỉnh Holm. \\
4.8 & Nguy cơ quá khớp (Overfitting) và độ tin cậy xác suất có được hiệu chuẩn không? & Có & Dừng sớm trên Val; hiệu chuẩn Temperature Scaling ($T \approx 1.35$), kiểm định độ bền khi khuyết phương thức và kiểm tra tổng quát hóa chuỗi thời gian ($N=600$, AUROC $0.9884$). \\
\bottomrule
\end{tabularx}
\end{table}

\paragraph{Kết luận tổng thể:} Theo nguyên tắc khắt khe nhất của PROBAST+AI, nếu bất kỳ một miền nào có nguy cơ sai lệch cao thì đánh giá tổng thể của nghiên cứu phải được xếp vào mức \textbf{High Risk of Bias / High Applicability Concern}. Việc thừa nhận trung thực điều này giúp nghiên cứu bảo vệ vững chắc đóng góp khoa học cốt lõi: đây là một công trình mẫu mực về \textit{tính nghiêm ngặt trong phương pháp luận học máy y tế}, tạo tiền đề kỹ thuật vững chắc trước khi tiến hành thu thập dữ liệu bệnh viện thực tế.
